% !TEX root = ../main.tex
% ---------------------------------------------------------------------
%  B3 -- Strukturierter Text (Skript Kap. 3, Anhang C, E; AS 2)
% ---------------------------------------------------------------------
\vlsec{cB3}{B3}{Strukturierter Text (ST)}
\label{sec:st}

\subsection{POU-Typen: Funktion, Funktionsblock, Programm}
\label{sec:pou}
\begin{ztab}{@{}lXXX@{}}
 & \R{FUNCTION} & \R{FUNCTION\_BLOCK} & \R{PROGRAM} \\ \midrule
Gedächtnis & \T{nein} & ja (pro Instanz) & ja \\
Instanz & keine & \T{muss} instanziiert werden & genau eine, global \\
Rückgabe & 1 Wert über Namen (+ \R{VAR\_OUTPUT}) & beliebig viele \R{VAR\_OUTPUT} & Variablen \\
Aufruf & \R{x := F(a,b);} & \R{fb(in:=a); y:=fb.out;} & \R{P(); P.var} \\
Typisch & Umrechnung, Formel & Zähler, Filter, Flanke, Regler, Automat & \R{MAIN}, an Task \\
\end{ztab}
\begin{formel}[Wann reicht eine FUNCTION? (MEP-Frage)]
Wenn das Ergebnis \emph{nur vom aktuellen Eingang} abhängt -- z.\,B.\
Sensor-Rohwert \ra bar. Keine Werte aus früheren Zyklen nötig \ra kein
Gedächtnis \ra Funktion genügt. Sobald $y(k-1)$, $u(k-1)$, ein Timer oder
ein Zustand gebraucht wird \ra \R{FUNCTION\_BLOCK}.
\end{formel}
\SW{Programm-Variablen} von aussen lesbar mit \R{DoSomeMath.rSumme}.
Debugger: Breakpoint F9, Schritt F10/F8. Eine Funktion zeigt bei jedem
Aufruf frisch initialisierte lokale Variablen.

\subsection{Deklaration}
\begin{lstlisting}
FUNCTION F_Name : REAL        // Name : Rückgabetyp
VAR_INPUT  a, b : REAL;   END_VAR   // Kopie hinein
VAR_OUTPUT err  : BOOL;   END_VAR   // zusätzliche Ausgänge
VAR_IN_OUT arr  : ARRAY[*] OF INT; END_VAR // Referenz
VAR        tmp  : REAL := 0.0; END_VAR   // lokal
F_Name := a + b;              // Rückgabe = Funktionsname
\end{lstlisting}
Form immer \R{Name : Typ := Startwert;}. Ohne Startwert: 0 / \R{FALSE} /
0.0. Kommentare \R{(* *)} oder \R{//}. Zuweisung \R{:=}, Vergleich \R{=},
ungleich \R{<>}.

\subsection{Datentypen und Zahlendarstellung}
\begin{ztab}{@{}llX@{}}
Typ & Bit & Bereich / Zweck \\ \midrule
\R{BOOL} & 1 & TRUE/FALSE \\
\R{SINT}/\R{USINT} & 8 & $-128..127$ / $0..255$ \\
\R{INT}/\R{UINT} & 16 & $-32768..32767$ / $0..65535$ \\
\R{DINT}/\R{UDINT} & 32 & $\pm2.1\cdot10^9$ / $0..4.3\cdot10^9$ \\
\R{REAL}/\R{LREAL} & 32/64 & IEEE 754, $\approx$7 / 15 Stellen \\
\R{BYTE}..\R{LWORD} & 8..64 & Bitmuster, nur Bitoperationen \\
\R{TIME} & 32 & \R{T\#5S}, \R{T\#500MS}, \R{T\#1M30S} \\
\end{ztab}
\SW{Literale} \R{181} $=$ \R{2\#1011\_0101} $=$ \R{16\#B5}. Eine
Hex-Ziffer $=$ 4 Bit.
\SW{Zweierkomplement} negativ: alle Bits invertieren, $+1$.
$-5$ als INT $=$ \R{16\#FFFB}; dasselbe Bitmuster als UINT $=65531$.
Ein negativer Wert mehr als positive (INT bis $-32768$).
\SW{REAL} $(-1)^S\cdot1.M\cdot2^{E-127}$; $6.5=$ \R{16\#40D00000}.
0.1 ist nicht exakt darstellbar \ra REAL nie mit \R{=} vergleichen,
sondern mit Toleranz.

\subsection{Operatoren und Kontrollstrukturen}
Priorität: Klammern, Funktionsaufruf, \R{NOT}, \R{* / MOD}, \R{+ -},
Vergleiche, \R{AND}, \R{XOR}, \R{OR}. Im Zweifel klammern.
\begin{lstlisting}
IF a THEN ... ELSIF b THEN ... ELSE ... END_IF
CASE state OF
  ST_A:        ...;
  ST_B, ST_C:  ...;       // mehrere Werte
  1..3:        ...;       // Bereich
ELSE ...;
END_CASE
FOR i := 0 TO 9 BY 1 DO ... END_FOR    // feste Anzahl
WHILE bed DO ... END_WHILE            // prüft VOR Durchlauf
REPEAT ... UNTIL bed END_REPEAT       // prüft NACH: min. 1x
EXIT;     // innerste Schleife verlassen
RETURN;   // Baustein verlassen
\end{lstlisting}
\begin{falle}[WHILE in einer SPS]
Die ganze Schleife läuft in \emph{einem} Zyklus. Eingänge ändern sich
währenddessen nie (PAE ist eingefroren) \ra hängt die Abbruchbedingung an
einem Eingang, entsteht eine Endlosschleife \ra Zykluszeit-Überschreitung,
Task blockiert. Der SPS-Zyklus \emph{ist} schon die Schleife \ra lieber
\R{FOR} mit fester Anzahl oder pro Zyklus einen Schritt.
\end{falle}

\subsection{Arrays, Strukturen, Aufzählungen}
\begin{lstlisting}
arBuf : ARRAY[0..99] OF REAL;          // Index ausserhalb = Laufzeitfehler
arM   : ARRAY[0..3, 0..3] OF INT;      // 2D
arInit: ARRAY[0..4] OF REAL := [1.0, 2.5, 3.0, 0.5, 4.0];

TYPE ST_Status :                       // unter DUTs anlegen
STRUCT
  bTaster : BOOL;
  rVolt   : REAL;
END_STRUCT
END_TYPE
arHist : ARRAY[0..99] OF ST_Status;    // Array aus Strukturen
arHist[k].rVolt := 2.5;                // Punktnotation

TYPE E_Monat : (JANUAR := 1, FEBRUAR, MAERZ); END_TYPE
eM : E_Monat := JANUAR;
\end{lstlisting}
\begin{falle}[ENUM beginnt nicht \emph{immer} bei 0]
Ohne Angabe zählt TwinCAT ab 0 aufwärts -- man darf aber Werte setzen
(\R{JANUAR := 1}). Ein ENUM ist nur Lesbarkeit, im Maschinencode steht
eine Zahl.
\end{falle}
\SW{Array zurückgeben} Rückgabetyp \R{ARRAY[0..99] OF REAL} geht (feste
Grösse). Variable Länge nur als \R{VAR\_IN\_OUT arr : ARRAY[*] OF INT;}
mit \R{LOWER\_BOUND(arr,1)} / \R{UPPER\_BOUND(arr,1)} (liefern DINT).

\subsection{Typkonvertierung}
Keine impliziten Umwandlungen: immer \R{QUELLE\_TO\_ZIEL(x)}.
\R{REAL\_TO\_INT} rundet kaufmännisch, \R{TRUNC} schneidet ab,
\R{INT\_TO\_REAL} verlustfrei, \R{WORD\_TO\_INT}/\R{BOOL\_TO\_INT} für
Bitmuster.
\begin{falle}[Ganzzahldivision]
\R{INT / INT} bleibt INT: \R{7/2 = 3}. Für Skalierungen zuerst
\R{INT\_TO\_REAL} und mit \R{10.0} statt \R{10} rechnen.
\end{falle}

\subsection{Bitoperationen und Statuswort}
\label{rez:statuswort}
\R{SHR(x,n)} $=$ $x/2^n$ (rechts schieben), \R{SHL(x,n)} $=$ $x\cdot2^n$.
Einzelbit direkt: \R{w.6}. Setzen mit \R{OR}, löschen mit \R{AND NOT},
toggeln mit \R{XOR}.
\begin{rezept}[Bitfeld aus einem WORD extrahieren]
Maske mit \R{AND} \ra mit \R{SHR} an Stelle 0 schieben \ra konvertieren:
\R{inLength := WORD\_TO\_INT(SHR(sw AND 16\#FF00, 8));}\\
Einzelbit: \R{bOverrun := sw.6;} oder
\R{(SHR(sw,6) AND 16\#0001) = 1}.
\end{rezept}
\begin{bsp}[AS 2: Shift und Statuswort]
\R{myWord := 16\#5D7A} (23930). \R{SHR(myWord,3)} $=$ \R{16\#0BAF} $=$
2991 (Rest .25 fällt weg).\\
\R{SHL(WORD\_TO\_DWORD(myWord),3)} $=$ \R{16\#0002EBD0} $=$ 191440. Ohne
DWORD-Wandlung fallen die oberen Bits raus (\R{16\#EBD0}).\\
\R{myDWord OR 16\#F0000000} $=$ \R{16\#F002EBD0} (oberste 4 Bit gesetzt).\\[2pt]
RS232-Statuswort:\\
{\scriptsize\begin{tabular}{@{}l*{8}{c}@{}}
Bit & 15..8 & 7 & 6 & 5 & 4 & 3 & 2 & 1/0 \\
 & InLength & -- & OR & FR & PA & BUF\_F & IA & RR/TA \\
\end{tabular}}\\[2pt]
\begin{tabularx}{\linewidth}{@{}lrX@{}}
\R{16\#1D06} & 29 & IA, RR \\
\R{16\#534C} & 83 & OR\_ERR, BUF\_F, IA \\
\R{16\#C705} & 199 & IA, TA \\
\end{tabularx}
\end{bsp}

\subsection{Standardbausteine: Timer und Flanken}
\label{sec:timer}
\begin{ztab}{@{}lX@{}}
FB & Verhalten (Eingang \R{IN}, Zeit \R{PT}, Ausgang \R{Q}, \R{ET}) \\ \midrule
\R{TON} & Einschaltverzögert: \R{Q} wird TRUE, wenn \R{IN} schon \R{PT} lang TRUE ist. \R{IN}=FALSE \ra Reset. \\
\R{TOF} & Ausschaltverzögert: \R{Q} folgt \R{IN} sofort auf TRUE, bleibt nach Abfall noch \R{PT} TRUE. \\
\R{TP} & Puls: steigende Flanke \ra \R{Q} genau \R{PT} lang TRUE, nicht retriggerbar. \\
\R{R\_TRIG}/\R{F\_TRIG} & \R{CLK} \ra \R{Q} einen Zyklus TRUE bei steigender/fallender Flanke. \\
\R{CTU}/\R{CTD}, \R{RS}/\R{SR} & Zähler, Flipflops (gleiches Muster) \\
\end{ztab}
\begin{lstlisting}
VAR tonBlink : TON; bLampe : BOOL; END_VAR
// selbstrücksetzender Takt: toggelt alle 167 ms (= 3 Hz)
tonBlink(IN := bAktiv AND NOT tonBlink.Q, PT := T#167MS);
IF tonBlink.Q THEN bLampe := NOT bLampe; END_IF
\end{lstlisting}
\begin{falle}[Timer-Fallen]
Den Timer \T{jeden Zyklus} aufrufen (nicht nur im IF) -- sonst läuft
\R{ET} nicht. Zum Neustarten muss \R{IN} einen Zyklus FALSE sein. Im
Automaten \R{IN := (state = ST\_X)} schreiben, dann startet er beim
Zustandseintritt und setzt sich beim Verlassen selbst zurück.
Blinkfrequenz $f$: halbe Periode $=\frac{1}{2f}$.
\end{falle}

\begin{bsp}[EdgeDetector mit Doppelklick (AS 2)]
\label{bsp:edge}
\begin{lstlisting}
FUNCTION_BLOCK EdgeDetector
VAR_INPUT  inputValue : BOOL; tDouble : TIME := T#400MS; END_VAR
VAR_OUTPUT risingEdge, fallingEdge, doubleClick : BOOL; END_VAR
VAR        lastValue : BOOL; tpWin : TP; END_VAR

risingEdge  := inputValue AND NOT lastValue;
fallingEdge := NOT inputValue AND lastValue;
lastValue   := inputValue;          // Wert für nächsten Zyklus
doubleClick := risingEdge AND tpWin.Q;  // 2. Flanke im Fenster
tpWin(IN := risingEdge AND NOT doubleClick, PT := tDouble);
\end{lstlisting}
Eine Instanz pro Taster (\R{edgeGelb}, \R{edgeBlau}) -- jede hat ihr
eigenes \R{lastValue}.
\end{bsp}
